\documentclass[11pt,letterpaper]{article}

\usepackage[margin=1in]{geometry}
\usepackage[T1]{fontenc}
\usepackage[utf8]{inputenc}
\usepackage{lmodern}
\usepackage{amsmath}
\usepackage{amssymb}
\usepackage{booktabs}
\usepackage{enumitem}
\usepackage[hidelinks]{hyperref}
\usepackage{setspace}

\setstretch{1.08}
\setlength{\parindent}{0pt}
\setlength{\parskip}{5pt}

\title{ECE 280 Lab 2 -- TA Solution and Troubleshooting Manual\\Signals and Systems: Signal Transformations in Hardware}
\author{}
\date{}

\begin{document}
\maketitle

Use this as a fast grading and troubleshooting guide for:
\begin{itemize}[leftmargin=1.2em]
    \item Lab 2 manual and report template submissions
    \item Hardware implementation of time shift, amplitude scaling, and time reversal
    \item MATLAB verification using measured Arduino-acquired data
\end{itemize}

\section{Quick Pass/Flag Menu}

\subsection*{Pass Immediately If All Are True}
\begin{itemize}[leftmargin=1.2em]
    \item Student shows required scope and MATLAB evidence for Steps 1.1 through 3.3.
    \item Required screenshots include student name/ID and are legible.
    \item Measured transformations are compared to predicted MATLAB transforms (not only raw captures).
    \item Report includes NRMSE values for $y_1,y_2,y_3,y_4$ and discusses largest error source.
    \item Final conclusions distinguish algorithmic intent from hardware non-idealities.
\end{itemize}

\subsection*{Flag for Review If Any Are True}
\begin{itemize}[leftmargin=1.2em]
    \item No evidence for one or more transformation steps (2.1--2.4 or 3.1--3.3).
    \item Shift/reversal claims are visual only, with no quantitative or annotated timing checks.
    \item MATLAB work appears synthetic-only (no linkage to acquired hardware datasets).
    \item Same screenshot reused across multiple required steps.
\end{itemize}

\section{Instructor Key: Core Expected Results}

Given base signal $x(t)$ and assigned parameters $t_0$ and $B$:
\[
\begin{aligned}
y_1(t) &= x(t-t_0) \\
y_2(t) &= B\,x(t) \\
y_3(t) &= x(-t) \\
y_4(t) &= B\,x(-t-t_0)
\end{aligned}
\]

Discrete-time mapping (length-$N$ sample vectors):
\[
\begin{aligned}
y_1[n] &\leftrightarrow \texttt{circshift(x,k)} \\
y_2[n] &\leftrightarrow \texttt{B*x} \\
y_3[n] &\leftrightarrow \texttt{fliplr(x)} \\
y_4[n] &\leftrightarrow \texttt{B*circshift(fliplr(x),k)}
\end{aligned}
\]
where $k = \lfloor t_0/T_s \rfloor$.

\textbf{TA note:} If students use non-periodic windows, circular shift boundary artifacts are expected and should be discussed, not treated as automatic failure.

\section{Pre-Lab Solution Guidance}

\subsection*{Required Conceptual Targets}
\begin{itemize}[leftmargin=1.2em]
    \item Correct derivations for $y_1,y_2,y_3,y_4$ from assigned base signal form.
    \item Correct sign convention for delay/advance and reversal.
    \item Correct mapping from continuous shift $t_0$ to integer sample shift $k$.
\end{itemize}

\subsection*{Accept/Revise Criteria}
\begin{itemize}[leftmargin=1.2em]
    \item Accept if algebra is correct and plotted/sketched outputs match transformed behavior.
    \item Revise if student confuses phase shift with amplitude scaling or claims reversal is same as negative amplitude.
\end{itemize}

\section{Measurement Targets and Acceptance Bands}

\subsection*{Step 1.1 Base Signal}
\begin{itemize}[leftmargin=1.2em]
    \item Frequency and amplitude should be close to TA-assigned targets.
    \item Typical acceptance band: within 5--10\% if setup is otherwise coherent.
\end{itemize}

\subsection*{Step 2.1 Amplitude Scaling}
\begin{itemize}[leftmargin=1.2em]
    \item Measured ratio $B_{meas}=A_{Ch1}/A_{Ch2}$ should track assigned $B$.
    \item Typical acceptance: ratio error under 10--15\% with clear annotation.
\end{itemize}

\subsection*{Step 2.2 Time Shift}
\begin{itemize}[leftmargin=1.2em]
    \item Cursor-measured delay should match target $t_0$ within sampling granularity.
    \item Expected resolution limit: approximately one sample period $T_s$.
\end{itemize}

\subsection*{Step 2.3 Time Reversal}
\begin{itemize}[leftmargin=1.2em]
    \item Reversed waveform should preserve amplitude envelope but reverse temporal order.
    \item Student should identify one distinct feature that appears mirrored in time.
\end{itemize}

\subsection*{Step 2.4 Combined Transform}
\begin{itemize}[leftmargin=1.2em]
    \item Combined output must simultaneously show reversal pattern, delay offset, and amplitude scaling.
\end{itemize}

\section{Numerical Verification Key (MATLAB)}

\subsection*{Expected Deliverables}
\begin{itemize}[leftmargin=1.2em]
    \item Five captured traces: $x,y_1,y_2,y_3,y_4$.
    \item Four overlays: measured vs predicted transform for each $y_i$.
    \item NRMSE table for all four transforms.
\end{itemize}

\subsection*{NRMSE Interpretation}
\[
\mathrm{NRMSE}(\%) = 100\times\frac{\|y_{meas}-y_{pred}\|_2}{\|y_{pred}\|_2}
\]

\begin{itemize}[leftmargin=1.2em]
    \item Lower NRMSE indicates better hardware-model agreement.
    \item Time shift and combined transforms commonly show larger error due to alignment and sampling jitter.
    \item High NRMSE is acceptable if causes are explicitly diagnosed with evidence.
\end{itemize}

\section{Code Key Checks (TA Spot-Check)}

\subsection*{Arduino Checks}
\begin{itemize}[leftmargin=1.2em]
    \item Gain variable updated from potentiometer read and applied as scalar multiplier.
    \item Delay implemented via circular buffer index offset (not by arbitrary blocking delay).
    \item Reversal implemented via reversed LUT indexing or equivalent index transform.
    \item Shared variables across ISR/main loop are marked \texttt{volatile} if interrupt-timed.
\end{itemize}

\subsection*{MATLAB Checks}
\begin{itemize}[leftmargin=1.2em]
    \item Predicted transforms generated from captured base $x$ (not hand-crafted replacement arrays).
    \item Uses correct operators (\texttt{circshift}, scalar multiply, \texttt{fliplr}).
    \item Axes labeled and overlay legends clear enough to audit claims.
\end{itemize}

\section{TA Checkoff Script (2--3 Minutes per Team)}
\begin{enumerate}[leftmargin=1.4em]
    \item Ask team to state assigned $t_0$ and $B$, then define each $y_i$ verbally.
    \item Verify one scope evidence image from Steps 2.1--2.4 with ID visible.
    \item Ask for one MATLAB overlay and corresponding NRMSE value.
    \item Ask: ``Which transform had the largest error and why?''
    \item Confirm final report includes all required evidence map items.
\end{enumerate}

\section{Troubleshooting Matrix}

\begin{table}[h!]
\centering
\caption{Common Failure Modes and Corrective Actions}
\begin{tabular}{p{0.24\textwidth} p{0.33\textwidth} p{0.35\textwidth}}
\toprule
\textbf{Observed Symptom} & \textbf{Likely Cause} & \textbf{TA Action} \\
\midrule
Amplitude scaling does not change with potentiometer & Pot wired incorrectly or wrong analog pin read & Verify 5V-GND on end terminals and wiper to configured analog input. \\
Measured delay is inconsistent across runs & Buffer index update bug or nonuniform sampling interval & Check circular-buffer indexing and enforce fixed sample timing. \\
Time reversal looks like inversion only & Index ordering incorrect; sign flip applied without reversing sample order & Require explicit reverse-index logic and compare with MATLAB \texttt{fliplr}. \\
Combined transform fails while individual steps pass & Wrong operation order (shift before/after reverse mismatch) & Re-derive $y_4=B\,x(-t-t_0)$ and align firmware/MATLAB operation order. \\
MATLAB overlays misaligned despite similar shapes & Time vector offset, trigger mismatch, or sample start mismatch & Apply alignment offset/cross-correlation before NRMSE computation. \\
Very noisy acquired traces & Ground/probe issues or serial acquisition instability & Recheck common ground, probe compensation, and acquisition window consistency. \\
\bottomrule
\end{tabular}
\end{table}

\section{Grading Menu Aligned to Lab 2 Expectations}
\begin{itemize}[leftmargin=1.2em]
    \item Pre-lab derivation quality: formulas, mappings, and sketches are correct.
    \item Hardware transformation evidence: complete and annotated scope captures.
    \item MATLAB verification: all required overlays and NRMSE metrics included.
    \item Technical analysis: strongest error sources identified and justified.
    \item Authenticity compliance: required screenshots include name/ID.
\end{itemize}

\subsection*{Suggested Deductions}
\begin{itemize}[leftmargin=1.2em]
    \item Missing required evidence image: \(-4\) to \(-10\), depending on impact.
    \item No quantitative error metric: \(-4\) to \(-8\).
    \item Incorrect transform mapping with no correction attempt: \(-5\) to \(-12\).
    \item Missing ID on required screenshots: \(-3\) to \(-10\).
\end{itemize}

\section{Minimum Passing Submission Standard}
A minimally acceptable submission should include:
\begin{itemize}[leftmargin=1.2em]
    \item Pre-lab derivations for all four transformed signals.
    \item Scope evidence for base signal and all four hardware transforms.
    \item MATLAB plots for captured signals and at least one valid measured-vs-predicted overlay.
    \item NRMSE values for transformation comparisons.
    \item Brief error analysis linked to measured evidence.
\end{itemize}

\end{document}
